\documentclass[11pt,a4paper]{article}
\usepackage[T1]{fontenc}
\usepackage[utf8]{inputenc}
\usepackage{lmodern}
\usepackage{microtype}
\usepackage[margin=27mm,headheight=14pt]{geometry}
\usepackage{amsmath,amssymb,amsthm,mathtools}
\usepackage{booktabs,array,longtable}
\usepackage{xcolor}
\usepackage{enumitem}
\usepackage{fancyhdr}
\usepackage{listings}
\usepackage{hyperref}
\usepackage{bookmark}
\definecolor{ink}{HTML}{16324F}
\definecolor{accent}{HTML}{19647E}
\definecolor{pale}{HTML}{EFF5F8}
\hypersetup{colorlinks=true,linkcolor=ink,citecolor=accent,urlcolor=accent,
 pdftitle={Small-prime congruences for A348410},
 pdfauthor={ChatGPT},pdfsubject={Binary supercongruences, Mobius transforms, and a sharp denominator theorem}}
\pagestyle{fancy}
\fancyhf{}
\fancyhead[L]{\small\textsc{Small-prime congruences for A348410}}
\fancyhead[R]{\small Research article}
\fancyfoot[C]{\thepage}
\renewcommand{\headrulewidth}{0.35pt}
\setlength{\parskip}{4pt}
\setlength{\parindent}{1.2em}
\setlength{\emergencystretch}{2em}
\setcounter{tocdepth}{2}
\numberwithin{equation}{section}
\newtheorem{theorem}{Theorem}[section]
\newtheorem{lemma}[theorem]{Lemma}
\newtheorem{proposition}[theorem]{Proposition}
\newtheorem{corollary}[theorem]{Corollary}
\theoremstyle{definition}
\newtheorem{definition}[theorem]{Definition}
\newtheorem{problem}[theorem]{Problem}
\newtheorem{example}[theorem]{Example}
\theoremstyle{remark}
\newtheorem{remark}[theorem]{Remark}
\newtheorem{conjecture}[theorem]{Conjecture}
\newcommand{\Z}{\mathbb Z}
\newcommand{\Q}{\mathbb Q}
\newcommand{\R}{\mathbb R}
\newcommand{\CT}{\operatorname{CT}}
\newcommand{\vp}{v_p}
\newcommand{\eps}{\varepsilon}
\newcommand{\coeff}[2]{[#1^{#2}]}
\newcommand{\A}{A_{\alpha,\beta}}
\newcommand{\B}{B_{\alpha,\beta}}
\newcommand{\dd}{\,\mathrm d}
\DeclareMathOperator{\Res}{Res}
\lstset{basicstyle=\ttfamily\small,columns=fullflexible,breaklines=true,
 frame=single,rulecolor=\color{ink!25},backgroundcolor=\color{pale},
 xleftmargin=0.5em,xrightmargin=0.5em,aboveskip=1em,belowskip=1em}

\begin{document}
\begin{titlepage}
\thispagestyle{empty}
\vspace*{11mm}
{\color{accent}\large\scshape Combinatorics and arithmetic of integer sequences\par}
\vspace{10mm}
{\color{ink}\LARGE\bfseries Completing the small-prime\par
congruences of OEIS A348410\par}
\vspace{6mm}
{\Large A binary defect formula and a sharp\par denominator-12 theorem\par}
\vspace{10mm}
{\large Research article prepared with ChatGPT\par}
{19 September 2026\par}
\vspace{9mm}
\hrule
\vspace{5mm}
\begin{abstract}
Let $a(n)=[x^n]((1-x)(1-x^2))^{-n}$. An OEIS comment of Peter Bala
conjectures the cubic congruence $a(mp^r)\equiv a(mp^{r-1})\pmod{p^{3r}}$
for primes $p\ge5$. A source audit located a prior elementary proof candidate
for that assertion and a two-parameter extension; we therefore do not present
the odd-prime result as a new resolution. Instead, we investigate the binary
case explicitly left untreated in that source.

For $\A(n)=[x^n](1-x)^{-\alpha n}(1+x)^{-\beta n}$, with arbitrary
integer parameters, we prove an explicit first-defect congruence modulo
$2^{3r-2}$. It gives a sharp universal modulus $2^{3r-3}$, improved to
$2^{3r-2}$ when $\alpha\beta$ is even. Combining this with a self-contained
odd-prime proof gives $12n^{-3}\sum_{d\mid n}\mu(d)a(n/d)\in\Z$ for every
$n\ge1$. The constant $12$ is the smallest possible common denominator.
We also obtain a primitive-necklace divisibility theorem, explain the failure
of an overbroad periodic-coefficient argument, derive algebraic generating
functions and two correction terms in the asymptotic expansion, and provide
exact reproducible checks. Mathematical completeness of the written proofs
and global priority are separate matters: no claim of exhaustive novelty
verification or external peer review is made.
\end{abstract}
\vfill
\noindent\small
\textbf{Archive contents.} This article in \LaTeX\ and PDF; three verification
programs; exact CSV tables; test reports; build instructions; a source and
priority audit; and a draft OEIS comment. No computational conjecture is used
as a premise of a theorem.
\end{titlepage}

\begingroup
\small
\setlength{\parskip}{0pt}
\tableofcontents
\endgroup
\clearpage

\section{The problem, the source boundary, and the results}
\label{sec:intro}

\subsection{Why the target changed during the investigation}
The sequence A348410 counts nonnegative integer solutions of
\begin{equation}\label{eq:baskets}
 n=\sum_{i=1}^n (u_i+2v_i).
\end{equation}
Equivalently, $n$ indistinguishable objects are distributed among $n$ labelled
white baskets and $n$ labelled black baskets, and the black occupancies are
even. We set $a(0)=1$. The OEIS entry records the coefficient formula
\begin{equation}\label{eq:an}
 a(n)=[x^n]\frac1{(1-x)^n(1-x^2)^n}
\end{equation}
and, in a comment dated 21 February 2022, Bala's conjecture
\begin{equation}\label{eq:bala}
 a(mp^r)-a(mp^{r-1})\in p^{3r}\Z
 \qquad(p\ge5\text{ prime},\ m,r\ge1).
\end{equation}
The live entry still labels this a conjecture in the version inspected for
this article \cite{oeis348410}.

That label is not a reliable certificate of present-day openness. In
particular, the 2026 articles of Niu and Prodinger already discuss the
algebraic generating function \cite{niu,prodinger}. More importantly for
this investigation, a public repository note, \emph{A cubic tower for a
two-parameter coefficient family}, contains an elementary proof candidate
for \eqref{eq:bala}, for its full integer-parameter analogue, and for a
$3^{3r-1}$ ternary bound \cite{prior}. Its parameter convention is
$[x^n]((1+x)^u(1-x)^v)^n$; ours corresponds to $(u,v)=(-\beta,-\alpha)$.
The full odd-prime parameter conjecture is also stated on A352373
\cite{oeis352373}. The repository source explicitly leaves the binary case
without an assertion.

The research question used here is therefore the following more precise
completion problem, not a claim that the odd-prime conjecture had no prior
proof.

\begin{problem}[Binary completion]\label{prob:binary}
For the integer-parameter family
\begin{equation}\label{eq:family}
 \A(n)=[x^n](1-x)^{-\alpha n}(1+x)^{-\beta n},
 \qquad \alpha,\beta\in\Z,
\end{equation}
does cubic growth of the divisibility exponent persist at $p=2$, with a
fixed loss independent of the level? Can one determine the first possible
binary defect and obtain an optimal common denominator for the cubic
M\"obius transform?
\end{problem}

The sources checked did not supply the binary formula or the denominator
conclusions below. This is a statement about the checked source boundary,
not an exhaustive claim about all mathematical literature. The odd-prime
result is credited to the prior proof candidate and reproved here to make
the new completion independent of an external unreviewed argument.

\subsection{Main theorems}
All binomial coefficients with an integer upper argument use the generalized
convention
\begin{equation}\label{eq:genbinomial}
 \binom{z}{k}=\frac{z(z-1)\cdots(z-k+1)}{k!},\qquad k\ge0.
\end{equation}
In particular, this is an integer when $z$ is a negative integer.

\begin{theorem}[First binary defect]\label{thm:binary}
Let $\alpha,\beta\in\Z$, let $m\ge1$ be odd, and let $r\ge2$. Then
\begin{align}\label{eq:binary-defect}
 &\A(m2^r)-\A(m2^{r-1})\notag\\
 &\quad\equiv
 2^{3r-3}\alpha\beta
 \binom{(\alpha+\beta+1)m-1}{m-1}
 \pmod{2^{3r-2}}.
\end{align}
Consequently the difference is always divisible by $2^{3r-3}$, and is
divisible by $2^{3r-2}$ whenever $\alpha\beta$ is even. If both parameters
are odd and $m=1$, then its valuation is exactly $3r-3$, at every $r\ge2$.
\end{theorem}

The first lift $r=1$ is deliberately not included in the displayed defect
formula. It always satisfies the elementary congruence modulo $2$ proved
below. The distinction prevents a common error: extrapolating a later-level
binary estimate to the first lift.

\begin{theorem}[All primes for the basket sequence]\label{thm:allprime}
For every prime $p$ and all $m,r\ge1$,
\begin{equation}\label{eq:all-prime}
 v_p\bigl(a(mp^r)-a(mp^{r-1})\bigr)
 \ge 3r-v_p(12).
\end{equation}
Thus the moduli are $p^{3r}$ for $p\ge5$, $3^{3r-1}$ for $p=3$, and
$2^{3r-2}$ for $p=2$. Each of these uniform losses is necessary.
\end{theorem}

For a zero difference, its valuation is interpreted as $+\infty$.

\begin{theorem}[Sharp common denominators]\label{thm:denom}
Define
\begin{equation}\label{eq:Bn}
 \B(n)=\frac1{n^3}\sum_{d\mid n}\mu(d)\A(n/d),\qquad
 B(n)=B_{2,1}(n).
\end{equation}
Then $24\B(n)\in\Z$ for all integer parameters and all $n\ge1$.
If $\alpha\beta$ is even, $12\B(n)\in\Z$. In particular,
\begin{equation}\label{eq:denom12}
 \boxed{\quad
 \frac{12}{n^3}\sum_{d\mid n}\mu(d)a(n/d)\in\Z_{\ge0}
 \quad(n\ge1).\quad}
\end{equation}
The constant $12$ is minimal for the single sequence $a(n)$. The constant
$24$ is minimal for the unrestricted integer-parameter family.
\end{theorem}

For the combinatorial interpretation, let $P_n$ be the number of primitive
cyclic words of length $n$ in the weighted alphabet of pairs $(u,v)$ of
nonnegative integers, where $(u,v)$ has weight $u+2v$, and the total word
weight is $n$. ``Primitive'' means that no shorter word is repeated to form
it. Cyclic rotations are identified.

\begin{corollary}[Primitive-orbit divisibility]\label{cor:necklace}
For every $n\ge1$,
\begin{equation}\label{eq:necklace-div}
 P_n=\frac1n\sum_{d\mid n}\mu(d)a(n/d),\qquad
 \frac{n^2}{\gcd(n^2,12)}\mid P_n.
\end{equation}
\end{corollary}

\subsection{What is proved, and what is not claimed}
The binary theorem and denominator theorem have complete elementary proofs
in Sections~\ref{sec:binary} and~\ref{sec:mobius}. The odd-prime proof in
Section~\ref{sec:odd} is included as a self-contained alternative to the
prior repository argument. The algebraic generating function and leading
asymptotic are not claimed as new: both are already documented in the OEIS
and related literature. The correction coefficients and the general
Gaussian-moment expression are derived explicitly here, without a priority
claim. A stronger exact valuation pattern at powers of two is labelled as
conjectural in Section~\ref{sec:frontiers}.

The proofs have been checked algebraically and by independent exact
computations, but have not been externally refereed or formalized in Lean.
The code is a reproducibility aid, not a substitute for the arguments.

\section{Counting formulas and exact coefficient computation}
\label{sec:coeff}

\subsection{A positive binomial sum}
A single white basket contributes $(1-x)^{-1}$ and a black basket contributes
$(1-x^2)^{-1}$. Multiplying over the labelled baskets proves
\eqref{eq:an}. Expanding the two factors separately gives, for $n\ge1$,
\begin{equation}\label{eq:positive-sum}
 a(n)=\sum_{k=0}^{\lfloor n/2\rfloor}
 \binom{2n-2k-1}{n-2k}\binom{n+k-1}{k}.
\end{equation}
The first terms are
\[
 1,\ 1,\ 5,\ 19,\ 85,\ 376,\ 1715,\ 7890,\ 36693,\ 171820,\ 809380.
\]
This positive sum gives an independent check on the faster coefficient
recurrence below. It also makes integrality and nonnegativity immediate for
the target sequence.

For arbitrary integer parameters, generalized binomial expansion shows that
\begin{equation}\label{eq:signed-sum}
 \A(n)=\sum_{k=0}^n(-1)^k
 \binom{\beta n+k-1}{k}
 \binom{\alpha n+n-k-1}{n-k}\qquad(n\ge1).
\end{equation}
This formula includes zero and negative parameters without any change of
convention. We put $\A(0)=1$.

\subsection{A two-step recurrence in the coefficient index}
Fix $n$ and write
\[
 F(x)=(1-x)^{-\alpha n}(1+x)^{-\beta n}
      =\sum_{k\ge0}f_kx^k.
\]
Logarithmic differentiation gives
\[
 (1-x^2)F'(x)=n\bigl((\alpha-\beta)+(\alpha+\beta)x\bigr)F(x).
\]
Consequently, with $f_{-1}=0$ and $f_0=1$,
\begin{equation}\label{eq:coefficient-recurrence}
 kf_k=n(\alpha-\beta)f_{k-1}
      +\bigl(n(\alpha+\beta)+k-2\bigr)f_{k-2},\qquad k\ge1.
\end{equation}
The desired output is $f_n$. This is a recurrence in $k$ \emph{with $n$
fixed}, not a recurrence relating successive diagonal values $\A(n)$.

The algorithm uses $O(n)$ exact arithmetic operations and two live integer
coefficients. The integers grow with $n$, so this is not a claim of linear
bit complexity. Every division by $k$ is checked for exactness in the
supplied implementation. Large datasets are computed without floating-point
arithmetic.

\subsection{Small-prime boundaries visible in the first terms}
For the target sequence,
\begin{align*}
 a(2)-a(1)&=4,\\
 a(3)-a(1)&=18=2\cdot3^2,\\
 a(4)-a(2)&=80=5\cdot2^4,\\
 a(5)-a(1)&=375=3\cdot5^3,\\
 a(7)-a(1)&=7889=23\cdot7^3.
\end{align*}
These examples already show that $p=2$ and $p=3$ cannot simply be inserted
into Bala's stated modulus. They also prove the necessity of the fixed
losses in Theorem~\ref{thm:allprime}: at $p=2$ use level $r=2$, and at
$p=3$ use level $r=1$. At $p=5$ there is no extra uniform power available.

\section{Local arithmetic and the odd-prime theorem}
\label{sec:odd}

This section proves the odd-prime result independently. It overlaps in scope
with the prior coefficient-framing theorem \cite{prior}; the direct
coefficient-valuation argument below avoids dependence on the general
rational-framing statement discussed in Section~\ref{sec:boundary}.

\subsection{The local ring and a coefficient estimate}
For a prime $p$, let
\[
 R_p=\Z_{(p)}=\{u/v\in\Q:p\nmid v\}.
\]
Congruence modulo $p^h$ in this ring means that the difference has valuation
at least $h$. A formal power-series congruence is always coefficientwise.
All exponentials and logarithms are formal over $\Q$; at a fixed degree only
finitely many terms contribute.

Set
\begin{equation}\label{eq:phi-log}
 \Phi(x)=(1-x)^{-\alpha}(1+x)^{-\beta},\qquad
 L(x)=\log\Phi(x)=\sum_{j\ge1}\frac{c_j}{j}x^j,
 \qquad c_j=\alpha+\beta(-1)^j.
\end{equation}

\begin{lemma}[Valuation of a coefficient]\label{lem:gcoef}
Let $M\ge1$ and write $g_j=[x^j]\Phi(x)^M$. Then $g_j\in\Z$, $g_0=1$,
and for $j\ge1$,
\begin{equation}\label{eq:gvaluation}
 v_p(g_j)\ge\max\{0,v_p(M)-v_p(j)\}.
\end{equation}
\end{lemma}
\begin{proof}
Integrality follows from generalized binomial expansion. Differentiating
$\Phi(x)^M$ and multiplying by $x$ yields
\[
 jg_j=M\sum_{i=1}^j c_i g_{j-i}.
\]
The sum on the right is an integer. Thus $v_p(g_j)\ge v_p(M)-v_p(j)$,
and integrality gives the additional lower bound zero.
\end{proof}

\subsection{A reciprocal-square cancellation}
For odd $p$, put
\[
 \eps_p=\begin{cases}1,&p=3,\\0,&p\ge5.\end{cases}
\]

\begin{lemma}[Reduced reciprocal squares]\label{lem:squares}
If $s>0$, $p\mid s$, and $t=v_p(s)$, then
\begin{equation}\label{eq:Hs}
 H_p(s)=\sum_{\substack{1\le j<s\\p\nmid j}}\frac1{j^2}
 \in p^{t-\eps_p}R_p.
\end{equation}
\end{lemma}
\begin{proof}
Reduce denominators modulo $p^t$ and partition the summation interval into
complete blocks of length $p^t$. It is enough to consider
$U_t=\sum_{u\in(\Z/p^t\Z)^\times}u^{-2}$.

For $p\ge5$, multiplication by $2$ permutes the units and gives
$U_t\equiv 2^{-2}U_t\pmod{p^t}$. Hence $3U_t\equiv0\pmod{p^t}$; since
$3$ is invertible, $U_t\equiv0$.

For $p=3$, induction proves $U_t\equiv0\pmod{3^{t-1}}$. The case $t=1$
has modulus $1$. For $t\ge2$, each unit modulo $3^{t-1}$ has three lifts,
so $U_t\equiv3U_{t-1}\pmod{3^{t-1}}$. The induction hypothesis finishes
the proof. The same bounds hold for the full block sum.
\end{proof}

For odd $p$, parity is preserved by multiplication by $p$, so $c_{pj}=c_j$.
The reduced logarithm
\begin{equation}\label{eq:Dp}
 D_p(x)=L(x)-\frac1pL(x^p)
       =\sum_{\substack{j\ge1\\p\nmid j}}\frac{c_j}{j}x^j
\end{equation}
belongs to $xR_p[[x]]$ and has no monomial whose exponent is divisible by
$p$.

\begin{lemma}[Quadratic-tail estimate]\label{lem:quadratic}
If $s>0$ is divisible by the odd prime $p$, then
\begin{equation}\label{eq:Tvaluation}
 T_p(s):=[x^s]D_p(x)^2\in p^{v_p(s)-\eps_p}R_p.
\end{equation}
\end{lemma}
\begin{proof}
Write $t=v_p(s)$. Since $p\mid s$, both $j$ and $s-j$ are units whenever
$p\nmid j$. Thus
\[
 T_p(s)=\sum_{\substack{1\le j<s\\p\nmid j}}
 \frac{c_jc_{s-j}}{j(s-j)}
 \equiv-\sum_{\substack{1\le j<s\\p\nmid j}}
 \frac{c_jc_{s-j}}{j^2}\pmod{p^t}.
\]
If $s$ is odd, the weight $c_jc_{s-j}=\alpha^2-\beta^2$ is constant, so
Lemma~\ref{lem:squares} applies directly.

If $s$ is even, the weights are $(\alpha+\beta)^2$ for even $j$ and
$(\alpha-\beta)^2$ for odd $j$. The weighted reciprocal-square sum is
exactly
\[
 (\alpha-\beta)^2H_p(s)+\alpha\beta H_p(s/2).
\]
Indeed, the excess even weight is $4\alpha\beta$, while replacing $j$ by
$2k$ divides the reciprocal square by $4$. Since $p$ is odd,
$v_p(s/2)=t$. Applying Lemma~\ref{lem:squares} to both terms proves the
claim.
\end{proof}

\subsection{Proof of the full odd-prime tower}
\begin{theorem}[Odd-prime coefficient tower]\label{thm:odd}
For every $\alpha,\beta\in\Z$, odd prime $p$, and $N>0$ divisible by $p$,
write $r=v_p(N)$. Then
\begin{equation}\label{eq:odd-tower}
 \A(N)-\A(N/p)\in p^{3r-\eps_p}\Z.
\end{equation}
\end{theorem}
\begin{proof}
Put $N=pM$. From \eqref{eq:Dp} one obtains the exact identity
\begin{equation}\label{eq:factor}
 \Phi(x)^N=\Phi(x^p)^M\exp\bigl(ND_p(x)\bigr).
\end{equation}
The factorial valuation bounds imply
\begin{equation}\label{eq:exp-trunc}
 \exp(ND_p)\equiv1+ND_p+\frac{N^2}{2}D_p^2
 \pmod{p^{3r-\eps_p}R_p[[x]]}.
\end{equation}
Here are explicit bounds, including the small exponential degrees. For
$p\ge5$ and $k=3$, $v_p(k!)=0$. For $k\ge4$,
$v_p(k!)\le(k-1)/4\le k-3$, so $kr-v_p(k!)\ge3r$.
For $p=3$ and $k\ge3$,
$v_3(k!)\le(k-1)/2\le k-2$, and $kr-v_3(k!)\ge3r-1$.
These estimates justify every discarded term of the formal exponential.

Extract degree $N$. The constant term of the exponential gives $\A(M)$.
The linear contribution is zero: $\Phi(x^p)^M$ is supported on exponents
divisible by $p$, whereas $D_p$ is supported away from them. Therefore
\begin{equation}\label{eq:quadratic-reduction}
 \A(N)-\A(M)\equiv
 \frac{N^2}{2}[x^N]\Phi(x^p)^M D_p(x)^2
 \pmod{p^{3r-\eps_p}}.
\end{equation}
Writing $g_j=[x^j]\Phi(x)^M$, the remaining coefficient is
\begin{equation}\label{eq:tail-sum}
 \sum_{j=0}^{M-1}g_j T_p\bigl(p(M-j)\bigr).
\end{equation}
We claim that each summand belongs to $p^{r-\eps_p}R_p$.
For $j=0$, this follows directly from Lemma~\ref{lem:quadratic} at $s=N$.
For $j>0$ with $v_p(j)<r-1$, the equality $v_p(M)=r-1$ gives
$v_p(M-j)=v_p(j)$. Lemmas~\ref{lem:gcoef} and~\ref{lem:quadratic} then
supply, respectively,
\[
 r-1-v_p(j)\quad\hbox{and}\quad 1+v_p(j)-\eps_p,
\]
whose sum is $r-\eps_p$.
For $v_p(j)\ge r-1$, we have $v_p(M-j)\ge r-1$, and the quadratic-tail
bound alone supplies $r-\eps_p$.

Finally, $v_p(N^2/2)=2r$ because $p$ is odd. Inserting the valuation of
\eqref{eq:tail-sum} into \eqref{eq:quadratic-reduction} proves
\eqref{eq:odd-tower}. The difference is an ordinary integer, so its local
divisibility is ordinary integer divisibility.
\end{proof}

For $N=mp^r$ with $p\mid m$, the exact valuation $v_p(N)$ is larger than
the displayed level. The theorem therefore implies the usual formulation
without any coprimality restriction on $m$.

\section{The binary defect: proof and sharpness}
\label{sec:binary}

At $p=2$, the equality $c_{pj}=c_j$ used in \eqref{eq:Dp} fails. The correct
replacement is a reflection argument. It separates the even and odd parts
before any exponential is truncated.

\subsection{An exact odd-harmonic normalization}
Work in $R_2=\Z_{(2)}$ and define
\begin{equation}\label{eq:U}
 U(x)=\sum_{\substack{j\ge1\\j\text{ odd}}}\frac{x^j}{j}
     =\frac12\log\frac{1+x}{1-x}.
\end{equation}
The factor $1/2$ in the logarithmic expression is harmless: the displayed
odd-power series has coefficients in $R_2$.

\begin{lemma}[Odd harmonic sums are exactly normalized]\label{lem:binary-H}
For every positive even $T$, let
\[
 h(T)=\sum_{\substack{1\le j<T\\j\text{ odd}}}\frac1j.
\]
Then
\begin{equation}\label{eq:h-exact}
 v_2(h(T))=2v_2(T)-2.
\end{equation}
In particular, $h(2j)/j^2$ is an odd element of $R_2$ for every $j\ge1$.
\end{lemma}
\begin{proof}
Write $T=2^e q$ with $q$ odd. First,
\begin{equation}\label{eq:unit2}
 \sum_{\substack{1\le u<2^e\\u\text{ odd}}}u^{-2}
 \equiv2^{e-1}\pmod{2^e}.
\end{equation}
For $e=1$ this is immediate. For $e\ge2$, the two lifts $u$ and
$u+2^{e-1}$ of an odd residue modulo $2^{e-1}$ have equal inverse squares
modulo $2^e$. Indeed their squares differ by a multiple of $2^e$.
Thus the sum at level $e$ is twice the sum at level $e-1$ modulo $2^e$,
which proves \eqref{eq:unit2} by induction.

Partitioning the interval $1\le j<T$ into $q$ blocks now shows that
$\sum_{j\text{ odd}}j^{-2}\equiv2^{e-1}\pmod{2^e}$. Pair $j$ with
$T-j$ to obtain the exact identity
\[
 h(T)=\frac T2
 \sum_{\substack{1\le j<T\\j\text{ odd}}}\frac1{j(T-j)}.
\]
Modulo $2^e$, the sum on the right is the negative of the reciprocal-square
sum. It consequently has valuation exactly $e-1$. The prefactor $T/2$
has valuation $e-1$ as well, proving \eqref{eq:h-exact}.
\end{proof}

The even power series $U(x)^2$ satisfies
\[
 [x^{2j}]U(x)^2
 =\sum_{\substack{1\le k<2j\\k\text{ odd}}}\frac1{k(2j-k)}
 =\frac{h(2j)}j.
\]
Define
\begin{equation}\label{eq:K}
 K(t)=\sum_{j\ge1}\frac{h(2j)}{j^2}t^j\in tR_2[[t]],
 \qquad V(t)=U(\sqrt t)^2.
\end{equation}
Here $U(\sqrt t)^2$ denotes the integer-exponent series obtained from the
even series $U(x)^2$ by replacing $x^2$ with $t$; no choice of an analytic
square root is involved. The identities we need are
\begin{equation}\label{eq:K-identities}
 V(t)=tK'(t),\qquad
 U(t)^2=\frac t2\frac{\mathrm d}{\mathrm dt}K(t^2),\qquad
 K(t)\equiv\frac{t}{1-t}\pmod2.
\end{equation}
The first two follow by coefficient comparison. The last follows from the
oddness assertion of Lemma~\ref{lem:binary-H}. This final residue, not just
a divisibility bound, will determine the leading binary digit.

\subsection{Separating the two coefficient problems}
Set
\[
 S=\alpha+\beta,\qquad \delta=\alpha-\beta.
\]
Let $M$ be even, put $s=v_2(M)\ge1$, and define
\[
 g(t)=(1-t)^{-SM}.
\]
The logarithmic identities
\begin{align*}
 \Phi(x)^{2M}&=(1-x^2)^{-SM}\exp\bigl(2\delta M U(x)\bigr),\\
 \Phi(t)^M&=g(t)\exp\bigl(-2\beta M U(t)\bigr)
\end{align*}
are exact. Only the even part of the first exponential contributes to degree
$2M$. Consequently
\begin{align}\label{eq:binary-factor}
 \A(2M)&=[t^M]g(t)\cosh\bigl(2\delta M U(\sqrt t)\bigr),\notag\\
 \A(M)&=[t^M]g(t)\exp\bigl(-2\beta M U(t)\bigr).
\end{align}
The hyperbolic cosine is interpreted as its even formal Taylor series,
so every exponent of $t$ in \eqref{eq:binary-factor} is an integer.

Since $v_2(k!)\le k-1$, for $k\ge3$ we have
\[
 v_2\left(\frac{(2M)^k}{k!}\right)\ge ks+1\ge3s+1.
\]
The same estimate applies to all even degrees $k\ge4$ in the cosine.
Thus modulo $2^{3s+1}$ we need only three coefficient expressions:
\begin{equation}\label{eq:I-def}
 I_1=[t^M]gU,\qquad I_2=[t^M]gU^2,\qquad J=[t^M]gV.
\end{equation}
Writing $\Delta=\A(2M)-\A(M)$, expansion of
\eqref{eq:binary-factor} gives
\begin{equation}\label{eq:delta-start}
 \Delta\equiv
 2\beta M I_1+2\delta^2M^2J-2\beta^2M^2I_2
 \pmod{2^{3s+1}}.
\end{equation}

\subsection{Reflection converts the linear term into a quadratic one}
The apparent obstacle in \eqref{eq:delta-start} is the linear term $I_1$.
Because $M$ is even and $U$ is odd,
\[
 [t^M]g(-t)U(t)=-I_1.
\]
Also
\[
 \frac{g(-t)}{g(t)}=
 \left(\frac{1-t}{1+t}\right)^{SM}
 =\exp\bigl(-2SMU(t)\bigr).
\]
Combining these identities yields
\begin{equation}\label{eq:reflection}
 I_1=\frac12[t^M]g(t)U(t)
       \left(1-\exp\bigl(-2SMU(t)\bigr)\right).
\end{equation}
The first exponential term in \eqref{eq:reflection} is $SMI_2$.
Every remaining term has a scalar coefficient with valuation at least
$2s$: for $k\ge2$,
\[
 v_2\left(\frac{(2SM)^k}{2k!}\right)
 \ge k(s+1)-1-v_2(k!)\ge ks\ge2s.
\]
Hence
\begin{equation}\label{eq:I1-reduction}
 I_1\equiv SMI_2\pmod{2^{2s}R_2}.
\end{equation}
Multiplication by $2\beta M$ makes the error divisible by $2^{3s+1}$.
Substituting in \eqref{eq:delta-start} and using $\beta S-\beta^2=\alpha\beta$
therefore gives
\begin{equation}\label{eq:delta-quadratic}
 \Delta\equiv2M^2\bigl(\alpha\beta I_2+\delta^2J\bigr)
 \pmod{2^{3s+1}}.
\end{equation}
This reflection is the step that replaces the odd-prime support cancellation.
It is valid because $M$ is even; this is exactly why the first binary lift
was separated from the theorem.

\subsection{Formal integration by parts supplies the third power}
For any $K_0(t)\in R_2[[t]]$, coefficient differentiation gives
\begin{equation}\label{eq:ibp}
 [t^M]g(t)tK_0'(t)
 =M[t^M]g(t)K_0(t)\left(1-\frac{St}{1-t}\right).
\end{equation}
Indeed, $[t^M]t(gK_0)'=M[t^M]gK_0$ and
$tg'/g=SMt/(1-t)$.

Using \eqref{eq:K-identities} in \eqref{eq:ibp} shows
\begin{equation}\label{eq:IJ-div}
 J\in MR_2,\qquad
 I_2=\frac M2 C,
\end{equation}
where
\begin{equation}\label{eq:C}
 C=[t^M](1-t)^{-SM}K(t^2)
       \left(1-\frac{St}{1-t}\right)\in R_2.
\end{equation}
Thus the $J$ term in \eqref{eq:delta-quadratic} is already zero modulo
$2^{3s+1}$, and we obtain the particularly useful congruence
\begin{equation}\label{eq:delta-C}
 \Delta\equiv\alpha\beta M^3 C\pmod{2^{3s+1}}.
\end{equation}
It immediately proves the universal valuation bound $3s$ and its one-bit
improvement when $\alpha\beta$ is even.

It remains to compute $C$ modulo $2$ when both parameters are odd.
In this case $S$ is even. Write $M=2^s m$ with $m$ odd.
Over $\mathbb F_2[[t]]$, generalized integer powers and the Frobenius identity
give
\[
 (1-t)^{-SM}\equiv(1-t^{2^s})^{-Sm},\qquad
 1-\frac{St}{1-t}\equiv1,
\]
while $K(t^2)\equiv t^2/(1-t^2)$. Since $2^s\ge2$, extraction of $t^M$
then gives
\begin{align}\label{eq:C-binomial}
 C&\equiv\sum_{j=0}^{m-1}\binom{Sm+j-1}{j}\notag\\
  & =\binom{(S+1)m-1}{m-1}\pmod2.
\end{align}
The last identity is the generalized hockey-stick identity, obtained by
multiplying $(1-z)^{-Sm}$ by $(1-z)^{-1}$ and extracting degree $m-1$.
It holds for every integer $S$, including negative values.

If $\alpha\beta$ is even, both sides of the desired leading-bit formula
vanish regardless of $C$. If $\alpha\beta$ is odd, insert
\eqref{eq:C-binomial} into \eqref{eq:delta-C}. Finally,
$M^3=2^{3s}m^3\equiv2^{3s}\pmod{2^{3s+1}}$.
Taking $s=r-1$ proves Theorem~\ref{thm:binary} in full. For $m=1$ the binomial
coefficient in \eqref{eq:binary-defect} equals $1$, so odd parameters give
an odd leading quotient and exact valuation $3r-3$.

\subsection{First lift, optimal losses, and the all-prime statement}
For any $\Phi\in1+x\Z[[x]]$,
\[
 \Phi(x)^2\equiv\Phi(x^2)\pmod2.
\]
Raising this to $m$ and extracting degree $2m$ proves
\begin{equation}\label{eq:firstlift}
 \A(2m)\equiv\A(m)\pmod2.
\end{equation}
For $(\alpha,\beta)=(2,1)$, the improved binary bound of
Theorem~\ref{thm:binary} is $3r-2$ for $r\ge2$, and
\eqref{eq:firstlift} supplies exactly the claimed lower bound $1$ at $r=1$.
If $m$ is even, apply the theorem at the larger exact binary level.
Together with Theorem~\ref{thm:odd}, this proves
Theorem~\ref{thm:allprime}. The examples in Section~\ref{sec:coeff} prove
its uniform sharpness.

There is a distinction between uniform sharpness and eventual sharpness.
For the target sequence, equality occurs at $r=2$, but later powers of two
show more divisibility. The theorem does not assert equality at every level.
For the unrestricted family, in contrast, odd parameters and $m=1$ attain
the general bound $3r-3$ at every level $r\ge2$.

\section{M\"obius transforms, common denominators, and cyclic words}
\label{sec:mobius}

\subsection{From a prime-power tower to a denominator theorem}
For any sequence $A(n)$, set
\[
 T_A(n)=\sum_{d\mid n}\mu(d)A(n/d).
\]
If $n=p^r n_0$ with $p\nmid n_0$, separating divisors according as they
contain $p$ gives the exact identity
\begin{equation}\label{eq:mobius-group}
 T_A(n)=\sum_{d\mid n_0}\mu(d)
 \left(A(p^r n_0/d)-A(p^{r-1}n_0/d)\right).
\end{equation}
Only squarefree divisors matter, so there are no terms containing a higher
power of $p$ in the divisor itself.

Apply \eqref{eq:mobius-group} to the basket sequence. Theorem~\ref{thm:allprime}
implies, for every prime $p\mid n$,
\[
 v_p(T_a(n))\ge3v_p(n)-v_p(12).
\]
Therefore $n^3\mid12T_a(n)$, proving the integrality part of
\eqref{eq:denom12}.

For the general family, the odd-prime bounds are unchanged and the binary
bound loses at most three powers. At level one the needed general estimate
is merely valuation at least zero. Consequently
\[
 n^3\mid24T_{\A}(n).
\]
If $\alpha\beta$ is even, the improved binary estimate and
\eqref{eq:firstlift} reduce $24$ to $12$. This proves all the integrality
assertions of Theorem~\ref{thm:denom}.

The constants cannot be reduced. For the target sequence,
\begin{equation}\label{eq:denom-sharp}
 B(3)=\frac{19-1}{27}=\frac23,
 \qquad B(4)=\frac{85-5}{64}=\frac54.
\end{equation}
Every integer clearing all denominators must thus be divisible by both $3$
and $4$. Their least common multiple is $12$. For the unrestricted family,
$(\alpha,\beta)=(1,1)$ gives
\[
 B_{1,1}(4)=\frac{10-2}{64}=\frac18,
\]
and the target example at $n=3$ forces a factor $3$, so a universal clearing
integer must be divisible by $24$. Even the single parameter pair $(1,3)$
requires $24$: its values at $n=3$ and $n=4$ are $-8/3$ and $59/8$.

\subsection{Primitive cyclic words and positivity}
Let $\mathcal A=\{(u,v):u,v\in\Z_{\ge0}\}$, with weight $u+2v$.
A word of length $n$ and total weight $n$ corresponds exactly to a solution
of \eqref{eq:baskets}. Although the alphabet is infinite, the weight
constraint permits only finitely many letters at each fixed $n$.

A word of length $n$ whose least period is $d\mid n$ is the repetition of
a primitive block of length $d$. Its total weight is $n$ precisely when the
block weight is $d$. A primitive cyclic orbit of such blocks has exactly
$d$ distinct linear representatives. Therefore
\begin{equation}\label{eq:period}
 a(n)=\sum_{d\mid n}dP_d.
\end{equation}
M\"obius inversion proves the first formula of \eqref{eq:necklace-div} and
also proves $P_n\ge0$. Since $B(n)=P_n/n^2$, nonnegativity in
Theorem~\ref{thm:denom} follows. Integrality of $12P_n/n^2$ is equivalent
to the second assertion of \eqref{eq:necklace-div}.

The first values illustrate that this is stronger than the ordinary
integrality of a necklace count:
\begin{center}
\begin{tabular}{r r r r}
\toprule
$n$ & $P_n$ & $B(n)=P_n/n^2$ & $12B(n)$\\
\midrule
1&1&$1$&12\\
2&2&$1/2$&6\\
3&6&$2/3$&8\\
4&20&$5/4$&15\\
5&75&$3$&36\\
6&282&$47/6$&94\\
7&1127&$23$&276\\
8&4576&$143/2$&858\\
\bottomrule
\end{tabular}
\end{center}
For example, $P_7=49\cdot23$, and $P_6$ is divisible by
$36/\gcd(36,12)=3$. The individual counts need not be multiples of $n^2$;
small-prime denominators are genuine.

\subsection{Lambert series and an Euler product}
Let $H(z)=\sum_{n\ge0}a(n)z^n$. Equation~\eqref{eq:period} gives the formal
Lambert expansion
\begin{equation}\label{eq:Lambert}
 H(z)-1=\sum_{n\ge1}nP_n\frac{z^n}{1-z^n}
       =\sum_{n\ge1}n^3B(n)\frac{z^n}{1-z^n}.
\end{equation}
All operations are valid formally because only finitely many divisors
contribute to a fixed coefficient. The denominator theorem can be stated
as follows: the cubic Lambert coefficients of $12(H-1)$ are nonnegative
integers, and $12$ is the least positive integer with this property.

A related Euler product will follow from the Lagrange series in the next
section:
\begin{equation}\label{eq:euler-product}
 \frac{y(z)}z=\prod_{n\ge1}(1-z^n)^{-P_n}.
\end{equation}
To verify it directly, take formal logarithms. The coefficient of $z^k$
on the right is $k^{-1}\sum_{n\mid k}nP_n=a(k)/k$, exactly the logarithm
on the left derived below.

\section{Algebraic generating functions}
\label{sec:gf}

The main generating-function identities are established facts for A348410
\cite{oeis348410,niu,prodinger}. We include a short derivation because it
connects the coefficient-power definition, the Euler product, and the
asymptotic analysis.

\subsection{Lagrange reversion for the full parameter family}
Let $y(z)$ be the unique series in $z\Z[[z]]$ satisfying
\begin{equation}\label{eq:reversion}
 y=z\Phi(y).
\end{equation}
Existence and uniqueness follow recursively from $\Phi(0)=1$. Write
$H_{\alpha,\beta}(z)=\sum_{n\ge0}\A(n)z^n$.
Lagrange inversion, applied to $\log\Phi(y)$, gives
\begin{align*}
 [z^n]\log\Phi(y(z))
 &=\frac1n[x^{n-1}]L'(x)\Phi(x)^n\\
 &=\frac1{n^2}[x^{n-1}]\frac{\mathrm d}{\mathrm dx}\Phi(x)^n
 =\frac{\A(n)}n.
\end{align*}
Since $y/z=\Phi(y)$, it follows that
\begin{equation}\label{eq:log-y}
 \log\frac{y(z)}z=\sum_{n\ge1}\frac{\A(n)}n z^n.
\end{equation}
Differentiating and using \eqref{eq:reversion} yields
\begin{equation}\label{eq:Hparam}
 H_{\alpha,\beta}(z)=\frac{zy'(z)}{y(z)}
 =\frac{1-y^2}
 {1-(\alpha-\beta)y-(\alpha+\beta+1)y^2}.
\end{equation}
For integer parameters, $\Phi$ is rational. Thus \eqref{eq:reversion}
becomes a polynomial equation after clearing denominators, and
$H_{\alpha,\beta}$ is algebraic.

\subsection{The target quartic}
For $(\alpha,\beta)=(2,1)$,
\begin{equation}\label{eq:target-param}
 z=y(1-y)^2(1+y),\qquad
 H(z)=\frac{1-y^2}{1-y-4y^2}.
\end{equation}
Eliminating $y$ gives
\begin{equation}\label{eq:quartic}
 \begin{split}
 z^2(4H-1)^4
 +zH(107H^3-107H^2+36H-4)\\
 +32H^3(1-H)=0.
 \end{split}
\end{equation}
Substitution of \eqref{eq:target-param} verifies the identity directly.
The partial derivative of the left side with respect to $H$, at
$(z,H)=(0,1)$, is $-32$. Thus the formal implicit-function theorem singles
out the unique branch with $H(0)=1$. In particular, no ambiguous choice of
an algebraic root is hidden in the generating-function description.

The symbolic verification program checks both the elimination and the
series identity through degree $20$. The coefficient recurrence
\eqref{eq:coefficient-recurrence} is used for that series check, so it does
not merely substitute one rearranged symbolic expression into another.

\section{A boundary test for periodic-coefficient arguments}
\label{sec:boundary}

\subsection{What the logarithmic method actually requires}
The odd-prime proof used two different ingredients: the equality $c_{pj}=c_j$
and the much stronger quadratic cancellation in Lemma~\ref{lem:quadratic}.
The first does not imply the second.

More generally, let $\Phi\in1+xR_p[[x]]$, and suppose
\[
 \log\Phi(x)=\sum_{j\ge1}c_jx^j/j,\qquad c_j\in R_p,\qquad c_{pj}=c_j.
\]
Define $A(n)=[x^n]\Phi(x)^n$ and $D_p$ as in \eqref{eq:Dp}. The
factorization and the vanishing linear term still hold. For odd $p$, all
terms of degree at least two in the exponential have valuation at least
$2v_p(N)$, so they yield a quadratic tower. A cubic tower requires an
additional estimate on the coefficients of $D_p^2$.

At the first layer and for $p\ge5$, this gives a particularly transparent
necessary and sufficient residue condition:
\begin{equation}\label{eq:obstruction}
 A(p)-A(1)\equiv
 -\frac{p^2}{2}\sum_{j=1}^{p-1}\frac{c_jc_{p-j}}{j^2}
 \pmod{p^3}.
\end{equation}
Indeed, the quadratic coefficient at degree $p$ is unaffected by positive
powers of $x^p$ from $\Phi(x^p)$, and
$1/(j(p-j))\equiv-1/j^2\pmod p$. This proves the formula without any
periodicity assumption beyond $c_{pj}=c_j$.

\subsection{A spacing-three counterexample}
Replace the even-occupancy restriction by occupancy divisible by $3$:
\begin{equation}\label{eq:bn-boundary}
 b(n)=[x^n](1-x)^{-n}(1-x^3)^{-n}.
\end{equation}
Its logarithmic coefficients are
$c_j=1+3\mathbf1_{3\mid j}$. They have period three and satisfy $c_{5j}=c_j$.
Nevertheless,
\begin{align*}
 b(1)&=1,\\
 b(5)&=\binom95+5\binom62=126+75=201,\\
 b(5)-b(1)&=200\not\equiv0\pmod{125}.
\end{align*}
The obstruction in \eqref{eq:obstruction} is exactly
\begin{equation}\label{eq:weighted-counterexample}
 \sum_{j=1}^4\frac{c_jc_{5-j}}{j^2}
 =1+\frac44+\frac49+\frac1{16}
 =\frac{361}{144}\equiv4\pmod5.
\end{equation}
Equation~\eqref{eq:obstruction} predicts the residue $75$ modulo $125$,
which agrees with $200\bmod125$.

Thus the success of period two is structural, not an automatic consequence
of rationality, periodicity, or Frobenius invariance. In the period-two
proof, the weighted square sum collapses to two ordinary reciprocal-square
sums. That special collapse is absent here.

\subsection{Scope of the literature correction}
M\"uller's arXiv version 1 of \emph{Wolstenholme Type Congruences and Framing
of Rational 2-Functions} states a general periodic weighted-harmonic bound
in Theorems 1.2 and 6.2 and a corresponding precise cubic framing bound in
Theorem 1.1 \cite{muller}. The value \eqref{eq:weighted-counterexample}
contradicts the printed harmonic bound at period $3$ and prime $5$.

The coefficient counterexample also fits the printed number-field
hypotheses, not merely a superficial formal analogy. Take
$K=\Q(\sqrt{-3})$, whose discriminant is $-3$, and
\[
 V(x)=\frac{x}{1-x}+\frac{3x^3}{1-x^3}.
\]
At every unramified rational prime $p\ne3$, its rational coefficients
satisfy $c_{mp^r}=c_{mp^{r-1}}$ exactly. Thus it is a rational 2-function
in the convention of that preprint. Formula (5.3) there identifies the
coefficient of its plus-framing with parameter one as
$[x^n]\exp(n\int V)=b(n)$, where the formal integration sends $x^j$ to
$x^j/j$. Prime $5$ is unramified and does not divide the period. The
precise bound in Theorem 1.1 would therefore give divisibility by $125$,
contrary to $b(5)-b(1)=200$.

A different period-four counterexample is already documented by the same
repository that supplied the prior odd-prime proof \cite{priorcounter}.
Accordingly, the discovery of a problem in the general framing statement is
not claimed as new here. The spacing-three example is a compact independent
boundary witness. The criticism is specifically of the printed uniform
bounds in version 1. A single exceptional prime does not refute a weaker
assertion that permits an unspecified finite set of additional exceptional
primes, and it says nothing by itself about unrelated framing results.
None of the proofs in this article uses the disputed statements.

\section{Asymptotic expansion of the basket sequence}
\label{sec:asymptotic}

The leading asymptotic is recorded in the OEIS entry, attributed to
Kotesovec \cite{oeis348410}. Here we derive it together with two explicit
correction terms and an all-order coefficient formula. The method is the
standard saddle-point approach to large powers of generating functions
\cite[Chapter VIII]{flajolet}.

\subsection{The saddle and exponential growth}
For this section only, put
\[
 \Phi(x)=\frac1{(1-x)(1-x^2)}.
\]
The positive saddle $\tau\in(0,1)$ satisfies
\[
 \tau\frac{\Phi'(\tau)}{\Phi(\tau)}=1.
\]
Since $x\Phi'/\Phi=x/(1-x)+2x^2/(1-x^2)$, this is equivalent to
$4\tau^2+\tau-1=0$. Hence
\begin{equation}\label{eq:saddle}
 \tau=\frac{\sqrt{17}-1}{8},\qquad
 \rho=\frac{\tau}{\Phi(\tau)}
 =\frac{51\sqrt{17}-107}{512},\qquad
 \lambda=\rho^{-1}=\frac{107+51\sqrt{17}}{64}.
\end{equation}
Numerically, $\lambda=4.9574747954140732506\ldots$.
Define the logarithmic cumulants
\begin{equation}\label{eq:kappa}
 \kappa_j=\left.
 \left(x\frac{\mathrm d}{\mathrm dx}\right)^j\log\Phi(x)
 \right|_{x=\tau}.
\end{equation}
Then $\kappa_1=1$ and
\begin{align}\label{eq:kappas}
 \kappa_2&=\frac{51-5\sqrt{17}}{16},&
 \kappa_3&=\frac{435-69\sqrt{17}}{32},\notag\\
 \kappa_4&=\frac{10731-2125\sqrt{17}}{128},&
 \kappa_5&=\frac{91545-20175\sqrt{17}}{128},\notag\\
 \kappa_6&=\frac{4008567-929185\sqrt{17}}{512}.&&
\end{align}

\begin{theorem}[Two correction terms]\label{thm:asymptotic}
As $n\to\infty$ through positive integers,
\begin{equation}\label{eq:asymptotic}
 a(n)=\frac{\lambda^n}{\sqrt{2\pi\kappa_2n}}
 \left(1+\frac{c_1}{n}+\frac{c_2}{n^2}+O(n^{-3})\right),
\end{equation}
where
\begin{equation}\label{eq:c1c2}
 c_1=\frac{-1683+95\sqrt{17}}{9248},\qquad
 c_2=\frac{119821-28605\sqrt{17}}{5030912}.
\end{equation}
The leading amplitude is equivalently
\begin{equation}\label{eq:amplitude}
 \frac1{\sqrt{2\pi\kappa_2}}=
 \frac{\sqrt{3+5/\sqrt{17}}}{4\sqrt\pi}.
\end{equation}
\end{theorem}

\subsection{Derivation, including control of the remainder}
Cauchy's coefficient formula on the circle $|x|=\tau$ gives
\begin{equation}\label{eq:cauchy}
 a(n)=\frac{\lambda^n}{2\pi}
 \int_{-\pi}^{\pi}
 \exp\left(n\left\{
 \log\frac{\Phi(\tau e^{i\theta})}{\Phi(\tau)}-i\theta
 \right\}\right)\dd\theta.
\end{equation}
The coefficients of $\Phi$ are positive and its support contains $0$ and
$1$. Equality in the triangle inequality
$|\Phi(\tau e^{i\theta})|\le\Phi(\tau)$ therefore requires
$e^{i\theta}=1$. On every closed subinterval bounded away from zero, the
integrand is exponentially smaller than the central contribution.

Near zero, choose the analytic logarithm agreeing with the real logarithm
at $\tau$. The saddle condition cancels the linear term, and the exponent
inside braces is
\begin{equation}\label{eq:local-exponent}
 -\frac{\kappa_2\theta^2}{2}
 +\sum_{k\ge3}\frac{\kappa_k(i\theta)^k}{k!}.
\end{equation}
In particular, its real part is at most $-c\theta^2$ for sufficiently small
$|\theta|$ and some $c>0$.

Set $u=\sqrt{n\kappa_2}\,\theta$. The central integral becomes a Gaussian
integral multiplied by
\[
 \exp\left(\sum_{k\ge3}
 \frac{\kappa_k i^k u^k}
 {k!\kappa_2^{k/2}}n^{1-k/2}\right).
\]
Expanding this expression and integrating term by term is justified at any
fixed order by analytic Taylor remainder estimates and Gaussian domination.
For specificity, split a small fixed central interval at
$|\theta|=n^{-2/5}$. Outside that shrinking interval the bound
$e^{-cn\theta^2}$ is exponentially small in $n^{1/5}$. Inside it, Taylor
expansion to sufficiently high finite degree has a uniform remainder;
rescaling and integrating its polynomial bounds against a slightly wider
Gaussian gives the stated power remainder. Odd powers have zero symmetric
Gaussian moment. Carrying the expansion two orders further than the last
retained even power bounds the relative remainder by $O(n^{-3})$ in
\eqref{eq:asymptotic}. The same argument works to every fixed order.

Using
\[
 \frac1{\sqrt{2\pi}}\int_{\R}u^{2\ell}e^{-u^2/2}\dd u=(2\ell-1)!!,
\]
one obtains
\begin{align}\label{eq:cumulant-corrections}
 c_1={}&\frac{\kappa_4}{8\kappa_2^2}
       -\frac{5\kappa_3^2}{24\kappa_2^3},\notag\\
 c_2={}&-\frac{\kappa_6}{48\kappa_2^3}
 +\frac{7\kappa_3\kappa_5}{48\kappa_2^4}
 +\frac{35\kappa_4^2}{384\kappa_2^4}\notag\\
 &-\frac{35\kappa_3^2\kappa_4}{64\kappa_2^5}
 +\frac{385\kappa_3^4}{1152\kappa_2^6}.
\end{align}
Substitution of \eqref{eq:kappas} simplifies to \eqref{eq:c1c2}; elementary
rationalization gives \eqref{eq:amplitude}. This proves the theorem.

\subsection{An explicit all-order coefficient formula}
The same calculation gives a finite expression for every correction
coefficient. For $j\ge0$, sum over nonnegative integers
$m_3,\ldots,m_{2j+2}$ satisfying
\[
 \sum_{k=3}^{2j+2}(k-2)m_k=2j,
 \qquad K=\sum_{k=3}^{2j+2}km_k.
\]
The integer $K$ is even. Then
\begin{equation}\label{eq:allorder}
 c_j=\sum
 (-1)^{K/2}(K-1)!!\,\kappa_2^{-K/2}
 \prod_{k=3}^{2j+2}\frac{\kappa_k^{m_k}}{m_k!(k!)^{m_k}}.
\end{equation}
For $j=0$, the empty product and $(-1)!!$ are both $1$, giving $c_0=1$.
The constraint comes from collecting the exponent of $n^{-1/2}$ in the
expanded exponential; the double factorial comes from its Gaussian moment.
This is a direct finite enumeration formula, not a recurrence with unknown
initial correction terms.

Numerically, $c_1\approx-0.1396307272454933$ and
$c_2\approx0.0003736029529450764$. The cancellation in $c_2$ is substantial,
so exact symbolic simplification is preferable to low-precision arithmetic.
The supplied numerical checks use 90 decimal digits. For illustration,
the signed relative error of the approximation through $c_2$ is about
$-9.6264\cdot10^{-9}$ at $n=100$ and $-6.1254\cdot10^{-10}$ at $n=250$.
Those numerical comparisons are checks, not the proof of the remainder.

\section{Reproducible computation and proof audit}
\label{sec:verification}

\subsection{What the programs compute}
The archive separates exact integer checks from optional symbolic and
high-precision checks. The program \texttt{verify.py} uses only the Python
standard library and implements \eqref{eq:coefficient-recurrence}; it
cross-checks against \eqref{eq:positive-sum} and \eqref{eq:signed-sum}.
The program \texttt{binary\_checks.py} tests the complete leading-defect
formula, the harmonic-unit lemma, and both common-denominator assertions.
The program \texttt{symbolic\_checks.py} uses SymPy and mpmath for
elimination, Gaussian coefficient enumeration, and asymptotic comparisons.

The recorded run produced the following exact checks. Each count is a finite
coverage statement, not an extrapolation to all indices.
\begin{center}
\small
\begin{tabular}{p{0.70\textwidth}r}
\toprule
Check & Cases\\
\midrule
Initial target terms &15\\
Target recurrence versus positive binomial sum &208\\
Signed parameter recurrence versus binomial sum &1,519\\
Target odd-prime congruences &420\\
Two-parameter odd-prime congruences &3,136\\
Odd-prime quadratic-tail tests &918\\
Original M\"obius/orbit checks &300\\
Original boundary witnesses &5\\
\midrule
Complete binary leading-defect formula &5,808\\
Binary first-lift congruence &1,936\\
Exact binary harmonic-unit normalization &400\\
Target denominator-12 and orbit divisibility &1,000\\
General denominator-24 / conditional denominator-12 &4,900\\
Binary sharpness and minimal-denominator witnesses &290\\
\bottomrule
\end{tabular}
\end{center}
All these checks passed. The binary leading-defect grid covers
$-5\le\alpha,\beta\le5$, odd $1\le m\le15$, and $2\le r\le7$, with
$m2^r\le2000$. The target odd-prime run uses odd primes at most $97$,
multipliers at most $12$ coprime to the prime, and indices at most $3000$.
The general odd-prime grid uses parameters from $-3$ to $3$, primes
$3,5,7,11,13$, and indices at most $350$.

The symbolic run additionally verifies the logarithmic derivative, the
quartic resultant, its parametrization, its power series through degree
$20$, the saddle point, the growth constant, the leading amplitude, and the
Gaussian formulas for $c_1$ and $c_2$. These tests used Python 3.13.5,
SymPy 1.14.0, and mpmath 1.3.0 in the recorded environment; the exact scripts
are not tied to those library versions.

\subsection{Reproduction commands and output data}
From the archive directory, run
\begin{lstlisting}
python verify.py
python binary_checks.py
python -m pip install -r requirements-optional.txt
python symbolic_checks.py
latexmk -pdf -interaction=nonstopmode -halt-on-error article.tex
\end{lstlisting}
Only the third and fourth commands need optional Python packages. No network
access is used by any of the checking programs themselves. The integer
checker accepts command-line coverage parameters; see \texttt{python
verify.py --help}.

The CSV files contain target terms through $1000$, the exact prime-power
valuation checks, the full binary parameter grid, primitive-orbit and
M\"obius values, and asymptotic comparisons. The reports preserve the actual
check counts. There are no randomized tests and no unrecorded dependence on
an external computer-algebra session.

\subsection{Independent points of possible failure}
Several interfaces deserve explicit scrutiny in any independent review.
First, a recurrence in the coefficient index must not be mistaken for a
recurrence in the diagonal index. Second, the odd-prime exponential has a
vanishing linear term, but the binary one does not; the latter is handled
by reflection, not by reusing the former proof. Third, the $1/2$ in formal
integration by parts at the binary prime is retained throughout and is
responsible for one of the losses. Fourth, the Möbius grouping is performed
at the exact prime-power exponent of $n$, not merely at a prime dividing
$n$. Finally, the finer binary valuation table is not used to justify any
of the uniform bounds.

The article uses only elementary formal series, local rational arithmetic,
coefficient extraction, and the stated saddle-point analysis. It does not
invoke the disputed periodic framing bound or any unverified asymptotic
recurrence from an OEIS comment.

\section{Further questions after the binary completion}
\label{sec:frontiers}

The uniform small-prime completion is proved, but it does not determine
all exceptional extra divisibility. For the basket sequence, the observed
valuations at $n=2^r$ are
\begin{center}
\begin{tabular}{r r r r r r r r r}
\toprule
$r$&1&2&3&4&5&6&7&8\\
\midrule
$v_2(a(2^r)-a(2^{r-1}))$&2&4&8&13&16&19&22&25\\
\bottomrule
\end{tabular}
\end{center}
The recorded exact computation continues this pattern through $r=15$.

\begin{conjecture}[A finer target valuation, not used above]\label{conj:finer}
For every $r\ge4$,
\[
 v_2\bigl(a(2^r)-a(2^{r-1})\bigr)=3r+1.
\]
\end{conjecture}

This is not a consequence of Theorem~\ref{thm:binary}, whose leading residue
vanishes for $(\alpha,\beta)=(2,1)$. Proving it would require following the
reflection calculation several binary digits further, including the next
residues of $K(t)$ and additional exponential terms. The statement for a
fixed odd multiplier is more varied: at $m=5$ the values from $r=4$ onward
in the checked range are $3r+3$, while at $m=7$ they are $3r+2$.

A second question is to characterize the periodic logarithmic coefficient
sequences for which the cubic tail estimate holds. Formula
\eqref{eq:obstruction} gives an immediate first-level obstruction, but
vanishing at the first level alone does not establish all higher levels.
The spacing-three example shows why a condition stronger than periodicity
and Frobenius invariance is needed. A useful classification would separate
identities forced by symmetry from congruences that depend on the prime.

A third question is to refine the common denominator for a fixed parameter
pair. The bounds $24$ and $12$ are optimal uniformly over the stated classes,
but special pairs can have smaller denominators. The exact first binary
residue already detects one such improvement. An analogous explicit
ternary residue could identify further parameter-dependent reductions.

These questions are distinct from the completed claims. In particular,
the remaining uncertainty about the finer valuation pattern does not limit
the proof of the universal denominator-12 theorem.

\section{Conclusion}

The initial OEIS label led to a useful but necessary correction of the
research target: a prior odd-prime proof candidate already existed. The
binary completion proved here uses a different mechanism, converting a
linear exponential contribution into a quadratic one by reflection and
then extracting an extra index factor by formal integration by parts.
An exact odd-harmonic normalization identifies the first possible binary
bit, not just a lower bound.

For A348410 the completed all-prime estimate is naturally encoded by the
single integer $12$. It clears every denominator of the cubic M\"obius
transform and is forced by the first few terms. The resulting primitive
cyclic-word divisibility gives a combinatorial expression of the same
arithmetic fact. The archive contains the full proofs and reproducible
checks; priority beyond the audited sources and external validation of the
new binary argument remain matters for further review.

\appendix
\section{A minimal exact evaluator}
The following code is the computational core; the full scripts add input
validation, independent formula checks, report generation, and CSV output.
\begin{lstlisting}[language=Python]
def coefficient_diagonal(n: int, alpha: int = 2,
                         beta: int = 1) -> int:
    if n < 0:
        raise ValueError("n must be nonnegative")
    f_minus_two, f_minus_one = 0, 1
    for k in range(1, n + 1):
        numerator = n * (alpha - beta) * f_minus_one
        numerator += (n * (alpha + beta) + k - 2) * f_minus_two
        current, remainder = divmod(numerator, k)
        if remainder:
            raise ArithmeticError("coefficient division was not exact")
        f_minus_two, f_minus_one = f_minus_one, current
    return f_minus_one
\end{lstlisting}
The running names refer to the two previously available coefficients at the
start of each iteration. At $n=0$ the function correctly returns $1$.

\section{Notation and dependency map}
\begin{center}
\small
\begin{tabular}{p{0.20\textwidth}p{0.73\textwidth}}
\toprule
Symbol & Meaning\\
\midrule
$\A(n)$ & Integer-parameter coefficient diagonal in \eqref{eq:family}.\\
$a(n)$ & The specialization $A_{2,1}(n)$, OEIS A348410.\\
$R_p$ & Rationals with denominator prime to $p$.\\
$D_p$ & Reduced odd-prime logarithm in \eqref{eq:Dp}.\\
$H_p(s)$ & Reduced reciprocal-square sum; not the generating function.\\
$h(T)$ & Sum of reciprocals of the positive odd integers less than $T$.\\
$U,K,V$ & Binary series in \eqref{eq:U} and \eqref{eq:K}.\\
$B(n)$ & Cubic M\"obius transform $n^{-3}\sum_{d\mid n}\mu(d)a(n/d)$.\\
$P_n$ & Primitive cyclic-word count, equal to $n^2B(n)$.\\
$H(z)$ & Ordinary generating function of $a(n)$.\\
$\kappa_j$ & Saddle cumulants in \eqref{eq:kappa}.\\
\bottomrule
\end{tabular}
\end{center}
The proof dependencies are short. Lemmas~\ref{lem:gcoef}--\ref{lem:quadratic}
imply Theorem~\ref{thm:odd}. Lemma~\ref{lem:binary-H}, reflection, and
\eqref{eq:ibp} imply Theorem~\ref{thm:binary}. These two theorems plus the
first binary lift imply Theorem~\ref{thm:allprime}. Grouping M\"obius divisors
then proves Theorem~\ref{thm:denom}; periodic word decomposition gives
Corollary~\ref{cor:necklace}. The generating-function and asymptotic sections
are independent checks and complementary structure, not premises of the
arithmetic proofs.

\begin{thebibliography}{99}
\bibitem{oeis348410}
OEIS Foundation Inc., \emph{A348410: Number of nonnegative integer solutions
to $n=\sum_{i=1}^n(a_i+b_i)$, with $b_i$ even}. Includes P.~Bala's
21 February 2022 coefficient formulas and supercongruence conjecture,
and V.~Kotesovec's leading asymptotic.
\url{https://oeis.org/A348410}. Inspected 19 September 2026.

\bibitem{oeis352373}
OEIS Foundation Inc., \emph{A352373}. Related parameterized
coefficient-power congruence.
\url{https://oeis.org/A352373}. Inspected 19 September 2026.

\bibitem{prior}
GitHub repository \texttt{rbajaj5/a183068-supercongruence},
\emph{A cubic tower for a two-parameter coefficient family},
\texttt{related-results/CoefficientFramingCubicTower.md}.
Public proof candidate; pinned repository snapshot
\texttt{3085b46fdbce} (full hash in the linked URL).
\href{https://github.com/rbajaj5/a183068-supercongruence/blob/3085b46fdbce945d156d2b5b8b9d1a66627b4375/related-results/CoefficientFramingCubicTower.md}{Permanent link to the inspected note}.
Inspected 19 September 2026. The source states no binary assertion.

\bibitem{priorcounter}
GitHub repository \texttt{rbajaj5/a183068-supercongruence},
\emph{Counterexample to the rational-framing theorem},
\texttt{related-results/RationalFramingCounterexample.md}, same pinned
snapshot as \cite{prior}.
\href{https://github.com/rbajaj5/a183068-supercongruence/blob/3085b46fdbce945d156d2b5b8b9d1a66627b4375/related-results/RationalFramingCounterexample.md}{Permanent link to the period-four counterexample}.
Inspected 19 September 2026; also cross-checked against the source's
result index and coefficient-framing discussion.

\bibitem{niu}
T.~Niu, \emph{An explicit algebraic generating function for OEIS A348410},
arXiv:2605.16553, 2026.
\url{https://arxiv.org/abs/2605.16553}.

\bibitem{prodinger}
H.~Prodinger, \emph{The generating function of A348410 in OEIS using the
diagonal method and another sequence (A001008) from OEIS},
arXiv:2605.21255, version 2, 23 May 2026.
\url{https://arxiv.org/abs/2605.21255}.

\bibitem{muller}
L.~Felipe M\"uller, \emph{Wolstenholme Type Congruences and Framing of Rational
2-Functions}, arXiv:2104.10754, version 1, April 2021.
\url{https://arxiv.org/abs/2104.10754}.
The criticism in Section~\ref{sec:boundary} concerns the precise printed
bounds of Theorems 1.1, 1.2, and 6.2 in this version, not all framing theory.

\bibitem{flajolet}
P.~Flajolet and R.~Sedgewick, \emph{Analytic Combinatorics}, Cambridge
University Press, 2009, especially Chapter VIII.
Authors' book site: \url{https://ac.cs.princeton.edu/home/}.
\end{thebibliography}
\end{document}
